The Compact Muon Solenoid (CMS) experiment at the Large Hadron Collider (LHC) is designed to explore particle physics at the TeV energy scale.
The electromagnetic calorimeter (ECAL) plays a central role in this program by providing precise energy measurements for electrons and photons produced in proton--proton collisions~\cite{Chatrchyan:1129810}.
In addition to its primary energy measurement capability, the ECAL can achieve excellent timing resolution, enabled by the fast scintillation response of lead tungstate (PbWO$_4$) crystals, the characteristics of the front-end electronics, and the high-frequency sampling of the readout system.

Future operation of the LHC, including Run~III and the High-Luminosity LHC (HL-LHC) era~\cite{Apollinari:2015wtw}, will significantly increase the instantaneous luminosity relative to previous runs.
While these conditions will allow the accumulation of unprecedented integrated luminosity, they will also substantially increase the number of overlapping proton--proton interactions per bunch crossing, commonly referred to as pileup (PU).
To mitigate the impact of pileup on event reconstruction, both the CMS and ATLAS experiments are pursuing the integration of precision timing information into their detectors~\cite{Butler:2019rpu,Morgenstern:2019xto}.
Within CMS, this strategy includes new dedicated timing detectors as well as upgrades to existing subsystems, such as the ECAL, with the goal of achieving time resolutions at the level of approximately 30\,ps in the HL-LHC era~\cite{Bruneliere:2006ra}.

The CMS ECAL is a hermetic homogeneous calorimeter composed of 75\,848 PbWO$_4$ scintillating crystals arranged in a barrel section and two endcap sections~\cite{cmstdr}.
The crystals are characterized by a high density, short radiation length, and small Moli\`ere radius, enabling a compact and highly granular detector geometry.
Scintillation light produced by electromagnetic showers is detected using avalanche photodiodes in the barrel and vacuum phototriodes in the endcaps~\cite{Bayatian:2006nff}.
The scintillation time profile of PbWO$_4$ is well matched to the LHC bunch crossing interval, with the majority of the emitted light collected within 25\,ns.

While the primary function of the ECAL is the measurement of electromagnetic energy, the achievable timing resolution has a direct impact on both energy reconstruction and physics sensitivity.
ECAL timing information is used to suppress backgrounds with broad or noncollision-related time distributions, including contributions from cosmic rays, beam halo muons, electronic noise, secondary interactions in detector material, and out-of-time proton--proton collisions~\cite{Bruneliere:2006ra}.
In addition, precise timing measurements significantly enhance sensitivity to searches for long-lived particles, where delayed decays can produce measurable shifts in the arrival time of photons and electrons at the calorimeter.

The reconstruction of ECAL time is sensitive to scintillation signals originating from pileup interactions, particularly those occurring in bunch crossings preceding or following the hard-scatter interaction of interest.
Such contributions, referred to as out-of-time pileup (OOT-PU), distort the reconstructed pulse shape and degrade the time resolution.
As luminosity increases, OOT-PU becomes a dominant limitation of existing timing reconstruction techniques.
Since ECAL timing is also used to reject backgrounds arising from OOT-PU, degradation in timing performance can adversely affect a wide range of physics measurements.

To address these challenges, it is essential to develop timing reconstruction methods that are robust against OOT-PU.
In this paper, a cross-correlation (CC) timing algorithm is introduced that exploits the pileup-aware ECAL energy reconstruction to mitigate the impact of out-of-time signals on the time measurement.
The performance of this algorithm is evaluated using data and simulation, and compared to the default Run~II ratio-based time reconstruction.
Improvements in single-crystal and object-level timing resolution are demonstrated, with direct relevance for Run~III operation and as a foundation for future precision timing applications at CMS.

The paper is organized as follows.
Section~\ref{sec:algo} introduces the ECAL time reconstruction algorithms, including the ratio-based reconstruction used in Run~II and the cross-correlation (CC) method developed to mitigate out-of-time pileup.
The timing performance, resolution methodology, and calibration strategy are presented in Section~\ref{sec:res}, together with results obtained in data and simulation.
The main conclusions are summarized in Section~\ref{sec:conclusion}.
